\documentclass[10pt,a4paper]{amsart}
\usepackage[margin=19mm]{geometry}
\usepackage{amsmath,amssymb,mathtools}
\usepackage[scr=rsfs,cal=euler]{mathalfa}
\usepackage{xfrac}
\usepackage[unicode,hidelinks,bookmarks=false]{hyperref}
\newcommand{\ie}{i.e.\ }
\newcommand{\cf}{cf.\ }
\newcommand{\resp}{respectively\ }
\newcommand{\Subsection}{\subsection}
\providecommand{\nobreakdash}{\nobreak\hbox{-}}
\setlength{\emergencystretch}{2em}
\multlinegap0pt
\newtheorem{theorem}{Theorem}
\newcommand{\R}{\mathbb{R}}
\newcommand{\Rpos}{\R_{\ge0}}
\renewcommand{\epsilon}{\ensuremath\varepsilon}
\newcommand{\fpm}{\mathbf{x}}
\newcommand{\fpmy}{\mathbf{y}}
\title{Counting thresholds for perfect matchings in hypergraphs}
\author{Strahinja Gvozdi\'c}
\date{Source: arXiv:2608.19345v1 (CC BY 4.0)}
\begin{document}
\maketitle

\noindent\emph{Source and licence.} Counting thresholds for perfect matchings in hypergraphs. Version arXiv:2608.19345v1 (CC BY 4.0), distributed by arXiv under the Creative Commons Attribution 4.0 International licence, http://creativecommons.org/licenses/by/4.0/, as recorded in the arXiv licence metadata for this exact version and shown on the abstract page https://arxiv.org/abs/2608.19345v1. Full licence notice: \texttt{licenses/arxiv-2608.19345-CC-BY-4.0.txt}.\par\smallskip

\noindent\textit{Editorial scope.} Counted unit: the article's theorem thm-ent-div (source0.tex lines 360--362). The article's complete original proof of it is delivered with it (source0.tex lines 364--398, one \texttt{proof} environment of 35 lines). No mathematics is written, repaired or re-derived: the statement and the proof are the article's own lines, in the article's own notation, and no retained line is substituted or re-flowed.  No further source-local context is needed: every cross-reference of every retained line resolves inside the two delivered spans.  Nothing was omitted from any retained span, so the package carries no omission record.  No retained line contains a citation, so the wrapper carries no bibliography and no bibliography entry.  No retained line contains a figure environment, an image-inclusion command or drawing code, so no image is delivered and the package declares no asset dependency.  Editorial formatting only, in addition: an AMS \texttt{amsart} preamble that restates the article's own declarations and macros used by the retained lines, namely the environments \texttt{theorem} on separate counters, together with the standard packages those lines need.  The retained environments print in this document as Theorem 1 in order of appearance, whereas the article prints the same items with the numbering of its own full document.  The complete native source of the article as distributed in the arXiv e-print for this exact version is bundled unchanged as \texttt{sources/208078/source0.tex}.
\begin{theorem}\label{thm-ent-div}
    For all integers $1\le d\le k-1$ such that $d\mid k$, it holds that $\beta_d(k)\le \beta_1(k/d)$.
\end{theorem}

\begin{proof}
    Fix $d$ and $k$ as in the statement of the theorem. Take arbitrary $\gamma,\epsilon>0$, and let $G$ be a $k$-uniform hypergraph on $n$ vertices with $\delta_d(G)\ge(\beta_1(k/d)+\gamma)\binom{n-d}{k-d}$, where $n$ is large enough. We need to show that
    $$h(G)\ge \frac{n}{k}\ln\left((1-\epsilon)\frac{k}{n}\,\frac{\binom{n}{d}}{\binom{k}{d}}\,\delta_d(G)\right).$$
    
    Let $H$ be an auxiliary $k/d$-uniform hypergraph on vertex set $\binom{V(G)}{d}$. Its edges are the sets $f\in\binom{V(H)}{k/d}$ such that $\bigcup f\in E(G)$, where $\bigcup f$ denotes $\bigcup_{U\in f} U$. Then for every $S\in \binom{V(G)}{d}$,
    \begin{align*}
        \deg^H_1S &=\left(\deg_d^GS\right)\frac{(k-d)!}{(k/d-1)!(d!)^{k/d-1}} \\
        &\ge (\beta_1(k/d)+\gamma)\binom{n-d}{k-d}\frac{(k-d)!}{(k/d-1)!(d!)^{k/d-1}} \\
        &\ge (\beta_1(k/d)+\gamma/2)\binom{\binom{n}{d}-1}{k/d-1},
    \end{align*}
    as long as $n$ is large enough (in terms of $\gamma$). Under the same hypothesis, now in terms of both $\gamma$ and $\epsilon$, it follows from the definition of $\beta_1(k/d)$ that there exists a fractional perfect matching $\fpmy:E(H)\to\Rpos$ of $H$ such that
    \begin{equation}\label{eq-ent-div}
        h(\fpmy)\ge \frac{\binom{n}{d}}{k/d}\ln\Big((1-\epsilon)\delta_1(H)\Big).
    \end{equation}

    For $e\in E(G)$ and $f\in E(H)$, write $e\sim f$ whenever $\bigcup f=e$, and let $a=1/\binom{n-1}{d-1}$. Define an edge-weighting $\fpm$ of $G$ by $\fpm_e:= a\sum_{e\sim f\in E(H)} \fpmy_f$. Then $\fpm$ is nonnegative. Moreover, for all $v\in V(G)$,
    \begin{align*}
        \sum_{e\ni v}\fpm_e &= a\sum_{e\ni v}\sum_{f\sim e} \fpmy_f = a\sum_{\substack{f\in E(H) \\ v\in\bigcup f}}\sum_{e\sim f}\fpmy \\
        &= a\sum_{\substack{f\in E(H) \\ v\in\bigcup f}}\fpmy_f = a\sum_{\substack{S\in\binom{V(G)}{d} \\ v\in S}} \sum_{S\in f\in E(H)} \fpmy_f \\
        &= a\sum_{\substack{S\in\binom{V(G)}{d} \\ v\in S}}1 = a\binom{n-1}{d-1} = 1,
    \end{align*}
    so $\fpm$ is a fractional perfect matching of $G$.
    
    Let $N=\frac{k!}{(k/d)!(d!)^{k/d}}$ be the number of edges $f\in E(H)$ such that $e\sim f$ for any fixed edge $e\in E(G)$. Then
    \begin{align*}
        h(G) \ge h(\fpm) &= \sum_{e\in E(G)}\Bigg( a\sum_{f\sim e}\fpmy_f\Bigg)\ln\frac{1}{aN }\frac{N }{\sum_{f\sim e}\fpmy_f} \\
        &= a\left(\ln\frac{1}{aN}\right)\sum_{f\in E(H)}\fpmy_f + aN\sum_{e\in E(G)} \frac{\sum_{f\sim e}\fpmy_f }{N}\ln\frac{N }{\sum_{f\sim e}\fpmy_f} \\
        &\ge \frac{n}{k}\ln\left(\frac{(k/d)!(d!)^{k/d}}{k!}\binom{n-1}{d-1}\right) + aN\sum_{e\in E(G)}\sum_{f\sim e}\frac{1}{N }\fpmy_f\ln\frac{1}{\fpmy_f} \\
        &= \frac{n}{k}\ln\left(\frac{(k/d)!(d!)^{k/d}}{k!}\binom{n-1}{d-1}\right) + ah(\fpmy) \\
        &\ge \frac{n}{k}\ln\left((1-\epsilon)\frac{(k/d)!(d!)^{k/d}}{k!}\binom{n-1}{d-1}\delta_1(H)\right) \\
        &= \frac{n }{k }\ln\left((1-\epsilon)\frac{(k/d)\,d!}{k\dots(k-d+1)}\binom{n-1}{d-1}\delta_d(G)\right) \\
        &= \frac{n}{k}\ln\left((1-\epsilon)\frac{k}{n}\,\frac{\binom{n}{d}}{\binom{k}{d}}\,\delta_d(G)\right),
    \end{align*}
    where the first inequality follows by Jensen's inequality applied to the concave function $x\mapsto x\ln(1/x)$, and the second one follows from (\ref{eq-ent-div}). This concludes the proof.
\end{proof}

\end{document}
